% 钱德拉塞卡拉·拉曼（综述）
% license CCBYSA3
% type Wiki

本文根据 CC-BY-SA 协议转载翻译自维基百科\href{https://en.wikipedia.org/wiki/C._V._Raman}{相关文章}

钱德拉塞卡拉·文卡塔·“C. V.”·拉曼爵士（Sir Chandrasekhara Venkata "C. V." Raman，/ˈrɑːmən/，\(^\text{[2]}\)音译“拉曼”；泰米尔语：சந்திரசேகர வெங்கட ராமன்，罗马化：Cantiracēkara Veṅkaṭa Rāmaṉ；1888年11月7日—1970年11月21日）是一位印度物理学家\(^\text{[3]}\)，以其在光散射领域的研究而闻名。\(^\text{[4]}\)他使用自己研制的分光仪，与学生K. S. 克里希南共同发现：当光线穿过透明物质时，被偏转的光会改变其波长。这种此前未知的光散射现象，他们称之为“修饰散射”，后来被命名为“拉曼效应”或“拉曼散射”。1930年，拉曼因这一发现获得诺贝尔物理学奖，成为第一位获得该奖项的亚洲人和非白人。\(^\text{[5]}\)

拉曼出生于泰米尔婆罗门家庭，自幼聪颖过人。他分别在11岁和13岁时，就完成了圣阿洛伊修斯英印高级中学的中学和高中课程。16岁时，他以优异成绩从马德拉斯大学附属的总督学院毕业，并以物理学荣誉学士的身份名列榜首。他的第一篇研究论文《光的衍射》发表于1906年，当时他仍是一名研究生。次年，他获得了硕士学位。19岁时，拉曼进入加尔各答的印度财政服务部门，担任副会计长。在那里，他结识了印度科学促进会（Indian Association for the Cultivation of Science，简称 IACS）——印度首个科研机构。该机构允许他独立开展研究，也正是在这里，拉曼在声学与光学领域作出了重要贡献。

1917年，拉曼受阿修托什·穆克吉任命，出任加尔各答大学下属拉贾巴扎尔理学院的首任帕利特物理学教授。在他首次前往欧洲的旅途中，亲眼目睹地中海的蓝色使他产生了强烈的好奇心——他认为当时流行的解释（即海洋的蓝色是由于天空中瑞利散射光的反射所致）是错误的。1926年，他创办了《印度物理学杂志》。1933年，他迁往班加罗尔，成为印度科学研究所的首位印度籍所长。同年，他还创立了印度科学院。1948年，他建立了拉曼研究所，并在此工作直至生命的最后时光。

拉曼效应于1928年2月28日被发现。印度政府将这一天定为“全国科学日”，每年举行纪念活动。
\subsection{早年生活与教育}
拉曼出生于英属印度马德拉斯省的蒂鲁吉拉帕利（Tiruchirappalli，今印度泰米尔纳德邦的蒂鲁吉拉帕利），父母均为泰米尔·艾耶尔婆罗门——父亲钱德拉塞卡拉·拉曼纳森·艾耶，\(^\text{[6][7]}\)母亲帕尔瓦蒂·阿玛尔。\(^\text{[8]}\)他是家中八个孩子中的第二个。\(^\text{[9]}\)他的父亲是一名中学教师，收入微薄。拉曼曾幽默地回忆说：“我出生时嘴里含着一把铜勺。因为在我出生时，父亲那‘辉煌’的薪水是每月十卢比！”\(^\text{[10]}\)1892年，全家迁往安得拉邦的维萨卡帕特南（当时称为Vizagapatam或Vizag），因为他的父亲被任命为A.V. 纳拉辛哈·饶夫人学院物理系教师。\(^\text{[11]}\)

拉曼在维萨卡帕特南的圣阿洛伊修斯英印高级中学接受教育。\(^\text{[12]}\)他在11岁时通过了中学毕业考试，13岁时以奖学金身份通过了艺术初级考试（First Examination in Arts，相当于今日的大学预科课程），\(^\text{[9][13]}\)并在两项考试中均名列安得拉邦学校委员会（今安得拉邦中等教育委员会）榜首。\(^\text{[14]}\)

1902年，拉曼进入马德拉斯（今钦奈）的总督学院就读，当时他的父亲已调任至该校教授数学与物理。\(^\text{[15]}\)1904年，他以优异成绩从马德拉斯大学获得文学学士（B.A.）学位，并在物理学与英语两门科目中均名列第一，荣获金牌奖。\(^\text{[14]}\)18岁时，仍为研究生的拉曼在英国期刊《哲学杂志》上发表了他的第一篇科学论文——《由矩形孔引起的不对称衍射带》（Unsymmetrical diffraction bands due to a rectangular aperture，1906年）。\(^\text{[16]}\)1907年，他以最高荣誉从同一所大学获得文学硕士（M.A.）学位。\(^\text{[17][18]}\)同年，他的第二篇论文发表于同一杂志，题为《液体的表面张力》。\(^\text{[19]}\)巧合的是，这篇论文与瑞利勋爵的一篇关于人耳对声音敏感度的论文同时刊登。\(^\text{[20]}\)由此，瑞利勋爵开始与拉曼通信，并礼貌地称呼他为“教授”。\(^\text{[14]}\)

拉曼的物理导师理查德·卢埃林·琼斯深知他的学术潜力，坚持建议他前往英国继续研究。\(^\text{[21]}\)琼斯为他安排了与吉法德上校的体检。然而，拉曼的健康状况欠佳，被认为是“体弱多病之人”。\(^\text{[22]}\)体检结果显示他无法适应英国严酷的气候。\(^\text{[11]}\)对此，拉曼后来回忆道：“吉法德检查了我，并证明如果我去英国，我会死于肺结核。”\(^\text{[23]}\)
\subsection{职业生涯}
拉曼的哥哥钱德拉塞卡拉·苏布拉马尼亚·艾耶早年已加入享有盛誉的印度政府部门——印度财政服务处（Indian Finance Service，今为印度审计与会计服务 Indian Audit and Accounts Service）。\(^\text{[24]}\)拉曼也步其后尘，于1907年2月的入职考试中名列第一，顺利取得任职资格。\(^\text{[25]}\)同年6月，他被派往加尔各答（今加尔各答）任副会计长。\(^\text{[11]}\)

在加尔各答期间，拉曼对印度科学促进会（Indian Association for the Cultivation of Science，简称 IACS）深感钦佩。\(^\text{[23]}\)该机构创立于1876年，是印度历史上第一所科学研究机构。他很快结识了几位重要人物：阿修托什·德伊，后来的长期合作伙伴；阿姆里塔·拉尔·瑟卡尔），IACS创办人马亨德拉拉尔·瑟卡尔之子、协会秘书；阿修托什·穆克吉，IACS执行委员、加尔各答大学校长。\(^\text{[14]}\)在他们的支持下，拉曼获准利用业余时间在IACS进行科学研究——哪怕是在“极其不寻常的时间”，\(^\text{[26]}\)正如他后来回忆的那样。\(^\text{[11]}\)当时，IACS尚未正式招募全职研究人员，也从未发表过科研论文。拉曼在1907年发表的论文《偏振光下的牛顿环》发表于《自然》杂志，成为该研究所的第一篇学术论文。\(^\text{[27]}\)这项成果极大地激发了IACS的科研热情，并促使其于1909年创办了《印度科学促进会公报》，其中拉曼是主要撰稿人。\(^\text{[14]}\)

1909年，拉曼被调任至英属缅甸仰光（今缅甸仰光），担任货币官。然而不久之后，他因父亲病逝而返回马德拉斯。父亲的去世及随后的丧礼使他在当地滞留了一整年。\(^\text{[28]}\)不久，他重新回到仰光任职，但在1910年又被调回印度马哈拉施特拉邦的那格浦尔。\(^\text{[29]}\)尚未在那格浦尔任职满一年，1911年他便被晋升为会计长，再次调往加尔各答。\(^\text{[28]}\)

从1915年起，加尔各答大学开始将研究生派至印度科学促进会（IACS），在拉曼的指导下进行科研工作。他的第一位学生是苏丹舒·库马尔·班纳吉（Sudhangsu Kumar Banerji），当时是加内什·普拉萨德指导下的博士研究生，后来成为印度气象局天文台总监。

自次年起，其他大学也纷纷效仿，如阿拉哈巴德大学、仰光大学、印多尔皇后学院、那格浦尔科学研究所、克里什纳斯学院以及马德拉斯大学等，\(^\text{[30]}\)都派学生前往IACS接受拉曼的指导。到1919年为止，拉曼已指导过十多名学生。\(^\text{[31]}\)同年，IACS秘书阿姆里塔·拉尔·瑟卡尔去世后，拉曼被授予两个荣誉职务：名誉教授与名誉秘书。\(^\text{[26]}\)拉曼将这一时期称为他人生的“黄金时代”。\(^\text{[32]}\)

加尔各答大学于1913年选定拉曼担任“帕利特物理学教授”，该职位由捐助人塔拉克纳特·帕利特爵士）设立。大学参议院于1914年1月30日正式任命，如会议记录所载：\\\\
“在1914年1月30日的参议院会议上，作出以下帕利特教授职位的任命：P. C. 雷博士与C. V. 拉曼先生。……每位教授的任命为长期制。教授须于年满六十岁时离任。”\(^\text{[14]}\)

在1914年之前，阿修托什·穆克吉曾邀请贾格迪什·钱德拉·玻色出任该职位，但玻色婉拒了邀请。\(^\text{[14]}\)作为第二人选，拉曼被任命为首任“帕利特物理学教授”，但由于第一次世界大战的爆发，他的上任被迫推迟。直到1917年，他加入加尔各答大学于1914年新设的拉贾巴扎尔理学院，才正式成为全职教授。\(^\text{[14]}\)拉曼在担任政府公务员十年后，勉强辞去职务——这一决定被称为一种“至高的牺牲”\(^\text{[26]}\)，因为教授的薪水几乎只有他原先收入的一半。不过，他的任职条件在加入大学时被明确写入报告，对他极为有利。报告中指出：\\\\
“C. V. 拉曼先生接受了塔拉克纳特·帕利特爵士物理学教授一职，条件是他无需被派往印度以外的地区……报告显示，C. V. 拉曼先生自1917年7月2日起正式就任帕利特物理学教授……拉曼先生已通知，他不必承担硕士或理学硕士课程的教学任务，以免影响其科研工作或指导高年级研究生进行研究。”\(^\text{[30]}\)

拉曼被任命为“帕利特物理学教授”一事，曾遭到加尔各答大学参议院部分成员的强烈反对，尤其是外国成员——理由是他既没有博士学位，也从未出国留学。作为回应，阿修托什·穆克吉特地安排加尔各答大学于1921年授予拉曼名誉理学博士学位，以示肯定。同年，他受邀前往牛津大学，在“大英帝国大学大会”上发表演讲。\(^\text{[34]}\)到那时，拉曼已在国际上享有盛誉，其接待方正是两位诺贝尔奖得主：J. J. 汤姆孙爵士与卢瑟福勋爵。\(^\text{[35]}\)1924年，拉曼当选为英国皇家学会院士后，穆克吉问他今后的目标，他笑答道：“当然是诺贝尔奖。”\(^\text{[26]}\)1926年，拉曼创办了《印度物理学杂志》，并担任首任主编。\(^\text{[36]}\)该杂志第二卷刊登了他著名的论文《一种新的辐射》，首次报道了“拉曼效应”的发现。\(^\text{[37][38]}\)

1932年，拉曼的“帕利特物理学教授”职位由德本德拉·莫汉·玻色接任。翌年，他被任命为位于班加罗尔的印度科学研究所（Indian Institute of Science，简称 IISc）所长，于1933年离开加尔各答。\(^\text{[39]}\)印度科学研究所的创立得到了多方支持：迈索尔国王克里希纳拉贾·瓦迪亚尔四世、实业家贾姆谢特吉·塔塔以及海得拉巴的尼扎姆——米尔·奥斯曼·阿里·汗爵士共同捐赠了土地与资金。1909年，印度总督明托勋爵批准建立该研究所，并由英国政府任命莫里斯·特拉弗斯为首任所长。\(^\text{[40]}\)拉曼成为第四任所长，也是第一位印度籍所长。在任期间，他聘用了G. N. 拉马钱德兰，后者后来成为著名的X射线晶体学家。1934年，拉曼创立了印度科学院，并创办了其学术期刊《印度科学院院刊》，该期刊后来分为三个分支刊物：《数学科学院刊》，《化学学报》，《地球系统科学学报》。\(^\text{[35]}\)大约在同一时期，拉曼早在1917年就提出设想的“加尔各答物理学会”也正式成立。\(^\text{[14]}\)

1943年，拉曼与他的前学生潘查帕凯萨·克里希纳穆尔蒂共同创办了“特拉凡哥尔化学制造有限公司”。\(^\text{[41]}\)该公司于1996年更名为“TCM有限公司”，是印度最早的有机与无机化学品制造企业之一。\(^\text{[42]}\)1947年，印度独立后新政府任命拉曼为首任“国家教授”。\(^\text{[43]}\)

拉曼于1948年从印度科学研究所退休，次年在班加罗尔创立了“拉曼研究所”。他担任所长，并一直在此工作，直至1970年去世。\(^\text{[43]}\)
\subsection{科学贡献}
\subsubsection{音乐声学}
\begin{figure}[ht]
\centering
\includegraphics[width=8cm]{./figures/3a7f4c6530b8aaba.png}
\caption{能级示意图：展示了参与拉曼信号的各个能级状态} \label{fig_QDLlm_1}
\end{figure}
拉曼的研究兴趣之一是音乐声音的科学基础。他受到赫尔曼·冯·亥姆霍兹所著《音的感觉》一书的启发，这本书是他加入印度科学促进会（IACS）后接触到的。\(^\text{[25]}\)在1916年至1921年间，拉曼大量发表了相关研究成果。他提出了基于速度叠加原理的弓弦乐器横向振动理论。他最早的研究之一是关于小提琴和大提琴中的“狼音”现象。\(^\text{[44][45]}\)此外，他还研究了各种小提琴及类似乐器的声学特性，包括印度传统弦乐器的声学行为。\(^\text{[46][47]}\)他甚至研究了水滴溅起的声音。\(^\text{[48]}\)拉曼还进行了他称之为“机械演奏小提琴实验”的研究。\(^\text{[49]}\)
\begin{figure}[ht]
\centering
\includegraphics[width=8cm]{./figures/8cb85b7f2b42ce80.png}
\caption{1930年诺贝尔奖颁奖典礼上的拉曼（左起第一位），与其他获奖者合影：C. V. 拉曼（物理学奖）、汉斯·费舍尔（化学奖）、卡尔·兰德施泰纳（医学奖）以及辛克莱·刘易斯（文学奖）。} \label{fig_QDLlm_2}
\end{figure}
拉曼还研究了印度鼓类乐器的独特性。\(^\text{[50]}\)他对塔布拉鼓与姆里当加姆鼓声音谐波特性的分析，是印度打击乐器声学的首批科学研究。\(^\text{[51]}\)他还撰写了一篇关于钢琴琴弦振动的批判性论文，探讨并修正了被称为考夫曼理论的观点。\(^\text{[52]}\)1921年短暂访问英国期间，拉曼抓住机会研究了伦敦圣保罗大教堂圆顶“回音廊”中的声音传播现象，该建筑能产生奇特的回声与声学效果。\(^\text{[53][54]}\)拉曼在声学方面的研究，无论在实验方法还是概念上，都为他后来的光学与量子力学研究奠定了重要基础。\(^\text{[55]}\)
\subsubsection{海水的蓝色}
在拓展光学研究的过程中，拉曼自1919年起开始研究光的散射现象。\(^\text{[56]}\)他在光学物理方面的第一个重要发现，便是对海水蓝色成因的解释。1921年9月，他从英国乘坐S.S. Narkunda号轮船返回印度途中，凝视地中海时，对海水的蓝色产生了浓厚兴趣。凭借随身携带的简单光学仪器——一台袖珍分光镜和一个尼科尔棱镜，他亲自对海水进行了观察与研究。\(^\text{[57]}\)当时，关于海水为何呈蓝色的多种假说中，\(^\text{[58][59]}\)最广为接受的是瑞利勋爵在1910年提出的解释：“深海那令人赞叹的深蓝色，与水本身的颜色无关，而只是天空蓝光在海面的反射。”\(^\text{[60]}\)瑞利对天空为何呈蓝色的解释是正确的——这一现象即如今所称的“瑞利散射”，即光在大气中被微粒散射与折射所致。\(^\text{[61][62]}\)因此，他关于海水蓝色来源于天空反射的观点在当时被普遍认为是正确的。然而，拉曼通过使用尼科尔棱镜消除了海面反射阳光的影响，直接观察海水后发现：海水在这种条件下显得比平时更加湛蓝。这一结果与瑞利的理论相矛盾，也使他确信，海水的蓝色并非仅仅来自天空的反射。\(^\text{[63]}\)

当S.S. Narkunda号轮船抵达孟买港（今孟买港）后，拉曼立即完成了一篇题为《海的颜色》的文章，并发表在1921年11月的《自然》杂志上。他在文中指出，瑞利的解释“通过一种简单的观察方式（使用尼科尔棱镜）就可以被质疑”。\(^\text{[63]}\)他写道：\\\\
“当通过一个放置在眼前的尼科尔棱镜观察海水，以消除表面反射时，可以看到阳光的光线穿入水中，并由于透视效应似乎在水的深处汇聚成一点。问题是：究竟是什么在散射光线，使得它的路径变得可见？一个值得考虑的有趣可能性是——这些散射光的粒子，至少部分，可能正是水分子本身。”\(^\text{[14]}\)
\begin{figure}[ht]
\centering
\includegraphics[width=6cm]{./figures/c373caa8734590ef.png}
\caption{拉曼《光的分子衍射》（Molecular Diffraction of Light，1922年）标题页} \label{fig_QDLlm_3}
\end{figure}
当拉曼抵达加尔各答后，他请来自仰光大学的学生K. R. 拉曼纳森在印度科学促进会（IACS）继续展开进一步研究。\(^\text{[64]}\)到1922年初，拉曼得出了自己的结论，并在《伦敦皇家学会学报》上发表了报告：\\\\
“本文拟提出一种完全不同的观点：在这一现象中，正如在天空呈色这一类似情形中一样，分子衍射决定了所观察到的亮度，并在很大程度上决定了其颜色。作为讨论的必要前提，本文将给出关于水中分子散射强度的理论计算与实验观察。”\(^\text{[65]}\)

正如拉曼所言，拉曼纳森于1923年发表了一篇详细的实验研究报告。\(^\text{[66]}\)随后，他于1924年对孟加拉湾的进一步研究提供了确凿的证据。\(^\text{[67]}\)如今人们已经知道，水本身的颜色主要源于其对光谱中红色与橙色长波光的选择性吸收——这种吸收来自水分子中O–H（氧与氢结合）伸缩振动模式的红外吸收泛频。\(^\text{[68]}\)
\subsubsection{拉曼效应}
\textbf{背景}

拉曼在光散射研究中的第二项重要发现，是一种全新的辐射现象——后来以他名字命名的“拉曼效应”。【69】在揭示了导致海水呈蓝色的光散射性质之后，拉曼将研究重点转向这一现象背后的原理。他在1923年的实验中发现：当阳光通过紫色滤光玻璃照射到某些液体和固体上时，除了入射光外，还会出现其他光线。这提示了可能存在新的散射形式。他的学生拉曼纳森当时认为这只是“荧光的微弱痕迹”。【14】1925年，新加入的研究助理K. S. 克里希南进一步指出，当光在液体中发生散射时，除了常见的偏振弹性散射外，还可能存在额外的一条散射谱线。\(^\text{[70]}\)他将这一现象称为“微弱的荧光”。\(^\text{[71]}\)然而，在接下来的两年中，对这一现象进行理论解释的尝试都未能取得成功。\(^\text{[72]}\)

这一研究的重大推动力来自“康普顿效应”的发现。1923年，圣路易斯华盛顿大学的阿瑟·康普顿发现了电磁波也可以用粒子来描述的证据。\(^\text{[73]}\)到1927年，这一现象已被科学界广泛接受，拉曼也认同这一观点。\(^\text{[74]}\)当年12月，康普顿获得诺贝尔物理学奖的消息传来时，拉曼兴奋地对克里希南说道：\\\\
“好消息……真是太好了！不过你看，克里希南——如果这对X射线成立，那么对可见光也一定成立。我一直这样认为。一定存在一个与康普顿效应对应的光学现象。我们必须去追寻它，而且我们正走在正确的道路上。它一定会被发现，而我们一定要拿下诺贝尔奖！”\(^\text{[75]}\)

然而，拉曼灵感的起源可以追溯得更远。正如康普顿后来回忆的那样：“也许正是那场多伦多的辩论，使他（拉曼）在两年后发现了拉曼效应。”\(^\text{[25]}\)所谓“多伦多辩论”，指的是1924年在多伦多举行的英国科学促进会会议上，关于“光量子是否存在”的讨论。会上，康普顿展示了自己的实验结果，而哈佛大学的威廉·杜安则提出相反的证据，主张光是波动。\(^\text{[76]}\)拉曼当时站在杜安一方，并直言道：“康普顿，你是个出色的辩手，但真理并不在你那一边。”\(^\text{[25]}\)

\textbf{散射实验}
\begin{figure}[ht]
\centering
\includegraphics[width=6cm]{./figures/ebf7f67e43fc5253.png}
\caption{拉曼与克里希南发表的早期苯的拉曼光谱\(^\text{[77]}\)} \label{fig_QDLlm_4}
\end{figure}
克里希南于1928年1月初开始了实验。\(^\text{[64]}\)1月7日，他发现无论使用哪种纯净液体，都会在可见光范围内产生偏振荧光。当拉曼看到这一结果时，他十分震惊——自己多年来竟从未观察到这种现象。\(^\text{[64]}\)当晚，他与克里希南将这一新现象命名为“修饰散射”，以区别于康普顿效应中的“未修饰散射”。2月16日，他们向《自然》杂志投稿了一篇题为《一种新的次级辐射》的论文，该文于3月31日发表。\(^\text{[78]}\)

1928年2月28日，他们首次成功获得了与入射光分离的“修饰散射”光谱。由于当时测定光波长的技术困难，他们最初依赖肉眼通过棱镜观察阳光所形成的颜色。拉曼为此发明了一种用于探测和测量电磁波的分光仪。\(^\text{[35][79]}\)谈到这项发明时，拉曼后来曾风趣地说：“当我获得诺贝尔奖时，我在设备上花的钱还不到200卢比。”\(^\text{[80]}\)不过，显然他在整个实验上的总花费远不止此数。\(^\text{[81]}\)从那一刻起，他们使用来自汞弧灯的单色光照射透明材料，并让透射光投射到分光仪上记录光谱。这样，散射光谱线便能够被精确测量与拍摄下来。\(^\text{[82][83]}\)

\textbf{公告}

同一天，拉曼向媒体正式宣布了这一发现。印度联合通讯社于次日，即2月29日报道了此事，标题为《辐射新理论：拉曼教授的发现》。\(^\text{[84]}\)报道写道：\\\\

“加尔各答大学的C. V. 拉曼教授作出了一项有望对物理学产生根本性意义的发现……
这一新现象所展现的特征，比康普顿教授用X射线所发现的还要令人震惊。其主要特征是：当物质被某一颜色的光激发时，其中的原子会发射出两种颜色的光，其中一种与激发光不同，且在光谱中位置更低。最令人惊讶的是，这种改变的颜色与所使用物质的种类完全无关。”\(^\text{[69]}\)

3月1日，《政治家报》以《原子对光的散射——新现象——加尔各答教授的发现》为标题，转载了这一新闻。\(^\text{[85]}\)拉曼于3月8日向《自然》杂志提交了一篇三段式简短报告，该文于4月21日发表。\(^\text{[86]}\)随后，他于3月22日提交了包含实际数据的论文，刊登于5月5日。\(^\text{[87]}\)3月16日，拉曼在班加罗尔召开的“南印度科学协会”会议上，以《一种新的辐射》为题，首次正式而详尽地介绍了这一发现。他的演讲发表于3月31日的《印度物理学杂志》。\(^\text{[37]}\)当天，他还将该论文的印刷版复印件一千份寄往世界各国的科学家，以便迅速传播这一重大发现。\(^\text{[38]}\)

\textbf{反响与结果}

最初，一些物理学家——尤其是法国和德国的学者——对这一发现的真实性持怀疑态度。耶拿弗里德里希·席勒大学的格奥尔格·尤斯曾致信慕尼黑大学的阿诺德·索末菲，询问道：“您认为拉曼关于液体中光学康普顿效应的研究可靠吗？……液体中散射谱线的清晰度让我感到怀疑。”索末菲随后尝试复现实验，但未能成功。\(^\text{[88]}\)1928年6月20日，柏林大学的彼得·普林斯海姆成功复现了拉曼的实验结果。他在接下来的论文中首次使用了“Ramaneffekt”（拉曼效应）和“Linien des Ramaneffekts”（拉曼谱线）等术语。\(^\text{[89][90]}\)英文版本“Raman effect”与“Raman lines”也随即被广泛采用。\(^\text{[91][23][92]}\)

除了作为一种全新的物理现象外，拉曼效应还是最早证明光具有量子性质的实验证据之一。1929年初，美国约翰斯·霍普金斯大学的罗伯特·W·伍德成为首位确认拉曼效应的美国科学家。\(^\text{[93]}\)他通过一系列实验验证后评论道：“在我看来，这一极其优美的发现，是拉曼对光散射现象长期而耐心研究的成果，也是量子理论最有力的实验证明之一。”\(^\text{[94][95]}\)以此为基础的“拉曼光谱学”由此诞生。1930年，英国皇家学会会长欧内斯特·卢瑟福在颁发休斯奖章时评价说，拉曼效应是“过去十年中实验物理领域最出色的三到四项发现之一”。\(^\text{[75]}\)

拉曼自信自己将获得诺贝尔物理学奖，但在1928年奖项授予欧文·理查森、1929年授予路易·德布罗意时，他颇感失望。他对1930年的获奖几乎笃定无疑——早在7月，他便预订了前往颁奖典礼的机票，尽管颁奖结果要到11月才公布。此后，他每天都仔细翻阅报纸，一旦没有看到诺贝尔奖的消息，便将报纸扔掉。\(^\text{[96]}\)最终，他如愿在当年获得了诺贝尔物理学奖。\(^\text{[97]}\)
\subsubsection{后期研究}
拉曼与瓦拉纳西的贝拿勒斯印度大学有着密切联系。他曾出席该校的奠基仪式\(^\text{[98]}\)，并在1916年2月5日至8日于该校举办的系列讲座中发表了关于数学与《物理学中的一些新方向》的演讲。\(^\text{[99]}\)此外，他还被聘为该校的常任访问教授。\(^\text{[100]}\)

1932年，拉曼与苏里·巴伽万塔姆共同确定了光子的自旋，进一步证实了光的量子性质。\(^\text{[101][92]}\)他与另一位学生纳根德拉·纳特合作，对声致光效应（即光被声波散射的现象）进行了理论解释，发表了一系列论文，形成了著名的“拉曼–纳特理论”。\(^\text{[102]}\)基于这一效应的调制器与开关系统，后来成为激光光学通信的重要技术基础。\(^\text{[103]}\)

他还进行了多项实验与理论研究，包括：对超声与高超声频率声波引起的光衍射的研究；\(^\text{[104][105]}\)以及X射线照射普通光下晶体所产生的红外振动效应的研究（1935年至1942年间陆续发表）。\(^\text{[106][107]}\)

1948年，拉曼通过研究晶体的光谱学行为，以一种全新的方式探讨了晶体动力学的基本问题。\(^\text{[108][109]}\)自1944年至1968年间，他持续研究了金刚石的结构与性质；\(^\text{[110][111]}\)20世纪50年代初，他又研究了多种具虹彩效应的矿物与物质的结构与光学性质，包括拉长石\(^\text{[112]}\)、珍珠长石\(^\text{[113]}\)、玛瑙\(^\text{[114]}\)、石英\(^\text{[115]}\)、蛋白石\(^\text{[116]}\)以及珍珠。\(^\text{[117]}\)此外，他还对胶体光学、电学与磁学各向异性进行了深入研究。\(^\text{[118][119]}\)进入20世纪60年代后，他的研究兴趣转向生物学领域，包括花朵的颜色机理与人类视觉的生理学特性等课题。\(^\text{[120][121][122]}\)
\subsection{个人生活}
1907年，拉曼与洛卡苏达丽·阿玛尔结婚。她是时任马德拉斯海关监督官S. 克里希纳斯瓦米·艾耶的女儿。\(^\text{[24]}\)婚礼的日期通常被记载为1907年5月6日，\(^\text{[123][124][125]}\)但据拉曼的侄孙女兼传记作者乌玛·帕拉梅斯瓦兰考证，实际日期应为1907年6月2日。\(^\text{[126][127]}\)这是一场由两人自行安排的婚姻，当时新娘年仅13岁。\(^\text{[128][43][129]}\)（不同资料对其年龄有所出入：若按其出生年份1892年来算，\(^\text{[23][124][125]}\)她应约15岁；但帕拉梅斯瓦兰确认她确为13岁，\(^\text{[130]}\)这一说法也与她的讣告相符——《当代科学》在她1980年5月22日逝世时记载其享年86岁。\(^\text{[131]}\)）她后来曾笑称，他们的婚姻与其说是因为她的音乐才华（两人初次相识时，她正在弹奏维纳琴 veena），不如说是“因为财政部门给已婚职员的额外津贴”。\(^\text{[43]}\)当时，这项“额外津贴”指的是每月为已婚官员增加150卢比的补贴。\(^\text{[24]}\)1907年夫妇二人搬到加尔各答后，曾一度被谣传改信基督教。原因是他们经常前往加尔各答的圣约翰教堂——洛卡苏达丽喜爱教堂的音乐，而拉曼则对教堂的声学效果充满兴趣。\(^\text{[42]}\)

他们育有两子：长子钱德拉塞卡尔·拉曼与次子文卡特拉曼·拉达克里希南，后者成为著名的无线电天文学家。由于思想分歧，拉曼后来与长子断绝了关系。\(^\text{[132]}\)拉曼的哥哥钱德拉塞卡拉·苏布拉马尼亚·艾耶的儿子——苏布拉马尼扬·钱德拉塞卡则在1983年获得了诺贝尔物理学奖。\(^\text{[133]}\)

拉曼一生中收集了大量具有独特光散射特性的石头、矿物与材料，这些藏品部分来自他的海外旅行，也有许多是他人赠送的礼物。\(^\text{[134]}\)他经常随身携带一台袖珍分光镜，用以研究各种样品。\(^\text{[135]}\)这些藏品连同他的分光仪，如今陈列于印度科学研究所（IISc）。\(^\text{[136][137]}\)

英国物理学家欧内斯特·卢瑟福在拉曼人生的多个关键时刻都起到了重要作用——他于1930年提名拉曼角逐诺贝尔物理学奖；同年，作为英国皇家学会会长，为他颁发了休斯奖章；1932年，又推荐他出任印度科学研究所所长一职。\(^\text{[11]}\)

拉曼对诺贝尔奖有着近乎执念的热情。他曾在加尔各答大学的一次演讲中说道：“我并不为被选为英国皇家学会院士（1924年）而感到受宠若惊。这只是一个小小的成就。若要说我真正渴望的，那就是诺贝尔奖。你们会看到，我将在五年内获得它。”\(^\text{[138]}\)他深知，如果自己真的获奖，就不能像往常那样等待诺贝尔委员会年底的公告，因为当时前往瑞典必须乘船，需要很长时间。\(^\text{[139]}\)于是，他满怀信心地在1930年7月预订了前往斯德哥尔摩的轮船票——两张，一张为自己，一张为妻子。\(^\text{[140]}\)在获得诺贝尔奖后不久的采访中，有人问他：“如果您更早发现拉曼效应，会发生什么？”他笑着回答：“那我可能就得与康普顿分享这项诺贝尔奖了，而我可不愿那样。我宁可独自获得整个奖项。”\(^\text{[141]}\)
\subsubsection{宗教观}
虽然拉曼很少谈论宗教，但他公开表示自己是一个不可知论者，【142】并反对被称为无神论者。【22】他的这种不可知论立场，主要受其父亲的影响——父亲崇尚赫伯特·斯宾塞、查尔斯·布拉德劳以及罗伯特·G·英格索尔等人的哲学思想。【143】他厌恶印度教的传统仪式，【144】但在家庭生活中并未完全摒弃这些习俗。【145】【146】此外，他也受到“不二论”哲学思想的影响。【147】他的标志性装束是头戴传统的印度头巾（pagri，下有发髻）与佩戴印度教圣线。尽管在南印度文化中并无戴头巾的习惯，他却笑着解释道：“哦，如果我不戴头巾，我的头都会胀大。你们都这么夸我，我得用头巾把自我包住。”【25】他甚至开玩笑地说，正是因为那顶头巾，他在首次访问英国时才获得了他人注意——尤其是J. J. 汤姆孙与卢瑟福的关注。【43】他曾在一次公开演讲中说道：\\\\
“没有天堂、没有天界、没有地狱、没有轮回、没有再生、也没有永生。唯一真实的事实是——人出生、活着、然后死亡。因此，一个人应当好好地度过自己的一生。”【148】

在一次与圣雄甘地（Mahatma Gandhi）和德国动物学家吉尔伯特·拉姆（Gilbert Rahm）的友好会谈中，谈话转向宗教话题。拉曼说道：\\\\
“我来回答你（拉姆）的问题。如果真的存在上帝，我们必须在宇宙中去寻找他；如果他不在那里，那就不值得去寻找……天文学与物理学的不断新发现，似乎正在不断揭示上帝的存在。”【22】

在生命的最后时刻，拉曼对妻子说：“我只相信人类的精神。”并嘱咐她：“我的葬礼要干净而简单——只需火化，不要那些繁琐的仪式。”【144】
\subsection{逝世}
1970年10月底，拉曼在实验室中突发心脏骤停并倒下。他被送往医院，医生在诊断后宣告，他的生命最多只剩下四个小时。【149】然而，拉曼顽强地又坚持了几天，并请求能回到研究所的花园中，在弟子与崇敬者们的陪伴下度过最后的时光。【150】

在去世前两天，他对一位昔日学生说道：“不要让科学院的期刊消亡。因为它们是衡量本国科学质量的敏感指标——它们能告诉人们，科学是否真正扎根在这片土地上。”【43】当晚，拉曼在卧室中与研究所管理委员会成员开会，讨论研究所未来的管理方向。【150】他还嘱咐妻子，在他去世后只需为他进行简朴的火葬，不要举行任何宗教仪式。1970年11月21日清晨，拉曼因自然原因辞世，享年82岁。【149】

拉曼逝世的消息传出后，时任印度总理英迪拉·甘地（Indira Gandhi）公开发表悼词：

> “全国、议会以及我们每一个人都为C. V. 拉曼博士的逝世而哀悼。
> 他是现代印度最伟大的科学家之一，也是我们国家悠久历史上最卓越的思想家之一。
> 他的头脑如同他所研究并阐释的钻石一般晶莹剔透。
> 他一生的事业在于揭示光的本质，而世界因他为科学所赢得的新知识而以各种方式向他致敬。”【151】
